\documentclass{beamer}

\usetheme {Madrid}

\input{my_beamer.sty}

\definecolor{myred}{RGB}{153, 18, 2}
\setbeamercolor{structure}{fg=myred}

\title{Determination of Drought Triggering Probability \\via Bayesian Network}
\author{Robert R. Tucci}
\institute{ICPAC, www.ar-tiste.xyz}

\begin{document}

\begin{frame}
\titlepage
\end{frame}

\begin{frame}{Links}


{\bf Repo} at
{\tiny \url{https://github.com/nishadhka/bn-ibf-e4drr} }

Includes
\begin{itemize}
\item python code (well documented with plenty of docstrings)
\item folder entitled {\tt jupyter\_notebooks} with 3 jupyter notebooks
\item folder entitled {\tt white-paper} with latex code and PDF for {\bf preliminary} white-paper 

{\tiny \url{https://github.com/nishadhka/bn-ibf-e4drr/blob/main/white-paper/anticipatory-action.pdf}}
\end{itemize}

I would appreciate it if you would take a look at the code and read the white paper. Comments from {\color{red} why did you do this instead of that}, or {\color{red} I think this part is wrong, or BS or too technical or unnecessary or obscure} are welcomed and I will attempt to address them.
\end{frame}

\begin{frame}{Brief review of Tim Palmer ECMWF Theory}

\begin{description}
\item[$I^*  =$] impact threshold 
   for drought, it equals minus an SPI (Standardized Precipitation Index)
   
\item[$\mu=$] ensemble member index

\item[$M=$] number of ensemble members, 51 for $ECMWF$

\item[$\indi()=$] indicator function. $\indi(a>b)$
equals 1 if $a>b$ and is zero otherwise

\item[$TD=$] triggering decision to do drought warning
\item[$P^*=$] threshold probability that triggers a  drought warning

\end{description}

$$
\underbrace{P(\rvI\geq I^*  )}_{
\text{probability that impact is greater than threshold $I^*$}}= \frac{1}{M} \sum_{\mu=1}^{M} \indi(I_{\mu} \geq I^*) 
$$



$$
TD=
\indi\left[P(\rvI\geq I^*  ) > P^*\right]
$$
\end{frame}

\begin{frame}{Problem to solve}

Find an automatic method of determining triggering
 probability $P^*$. 
 \begin{itemize}

 \item The new method should be better than Nishadh's method of doing this for Karamoja, The latter is based solely on FAR/HR/AUROC. 
 
 \item The new method should take
 into account {\bf antecedents} like population density, road density, historical catastrophes, soil moisture, etc.
 \end{itemize}

\end{frame}

\begin{frame}{Bnet for determining triggering 
probability $P^*$}

\begin{itemize}
\item {\bf bnet (Bayesian network)}, discussed in my free,
open source book 
{\tiny \url{https://github.com/rrtucci/Bayesuvius/raw/master/main.pdf}}
\item I do bnet calculations using {\bf Pyagrum} software at {\tiny \url{https://pyagrum.readthedocs.io/en/}} combined with a small amount of auxiliary 
code that I wrote.
\end{itemize}


$$
\xymatrix{
\rvX^1 \ar[dr]
&\rvX^2 \ar[d]
&\rvX^3\ar[dl]
\\
\rvP(\rvI> I^*)\ar[d] 
& \rvP^*\ar[dl]
&\rv{\eps}\ar[l]
\\
\rv{TD}
\ar[d]\ar[dr]\ar[drr]
\\
\ul{FAR} 
&\ul{HR}
&\ul{AUROC}
}
\begin{array}{l}
\text{$\rvX^1, \rvX^2, \rvX^3$ are antecedents}
\\
\text{$\rv{\eps}$ is noise}
\end{array}
$$


\end{frame}

\begin{frame}{Much Still to Do}

{\color{red} Need Help}

\begin{itemize}
\item {\bf Need real data to train bnet. 
} So far, I've only used synthetic  data to train it. 

\item Need to do a lot of work {\bf fine tuning} and
{\bf integrating} 
this software with full corpus of DAA 
(Drought Anticipatory Action) software

\item Need to do {\bf modelling of 
antecedent  probabilities}. This requires real data, which I haven't gotten yet.

\end{itemize}

\end{frame}

\end{document}